% GENERATED FILE. Edit canonical lessons, then run npm run generate:books.
% canonical-source-hash: fnv1a64:e8c94078c2363ffb
% canonical-lessons: KA-C126-nimbe, KA-C126-draksi, KA-C126-halasu, KA-C126-menasu, KA-C126-jirige, KA-R126-first-pass-the-house-and-things-you-carry, KA-R126-second-pass-clothes-and-food

\chapter{Lemon, Grapes, Jackfruit, Pepper, Cumin}
\label{ch:ka-lemon-grapes-jackfruit-pepper-cumin}
\hlchaptermodality{126}

\begin{chapteropening}
\textbf{By the end of this chapter:} \emph{I can say a lemon, grapes, a jackfruit, pepper and cumin.}

Complete the last lesson of chapter 126: I can say a lemon, grapes, a jackfruit, pepper and cumin.
\end{chapteropening}

\section[\texorpdfstring{\kn{ನಿಂಬೆ} (nimbe)}{ನಿಂಬೆ (nimbe)}]{\kn{ನಿಂಬೆ} (nimbe) — a lemon}

\label{lesson:KA-C126-nimbe}



\begin{quote}
Before the new one: say the Kannada for upma, then the Kannada for finger millet.
\end{quote}

\subsection*{You'll want to know: \kn{ನಿಂಬೆ}}
\textbf{\kn{ನಿಂಬೆ}} — \emph{nimbe} — \textquotedblleft{}a lemon\textquotedblright{}.

\begin{grammarlens}[title={What you've built}]
One more word for this chapter.
\end{grammarlens}

\subsection*{Your turn}
\begin{itemize}
\raggedright
  \item \emph{Say it:} \emph{nimbe}
  \item \emph{Say it:} \emph{nimbe}, once more
  \item \emph{Say it:} say \emph{nimbe}
  \item \emph{Recall:} say the Kannada for upma, then the Kannada for finger millet, then say \emph{nimbe} again
\end{itemize}

\begin{tcolorbox}[breakable,colback=teal!4,colframe=teal!35!black,title={Before you move on}]
What does \kn{ನಿಂಬೆ} mean? (\textquotedblleft{}A lemon\textquotedblright{}.) Say it once more.
\end{tcolorbox}
\section[\texorpdfstring{\kn{ದ್ರಾಕ್ಷಿ} (drākṣi)}{ದ್ರಾಕ್ಷಿ (drākṣi)}]{\kn{ದ್ರಾಕ್ಷಿ} (drākṣi) — grapes}

\label{lesson:KA-C126-draksi}



\begin{quote}
Before the new one: say the Kannada for finger millet, then the Kannada for a lemon.
\end{quote}

\subsection*{You'll want to know: \kn{ದ್ರಾಕ್ಷಿ}}
\textbf{\kn{ದ್ರಾಕ್ಷಿ}} — \emph{drākṣi} — \textquotedblleft{}grapes\textquotedblright{}.

\begin{grammarlens}[title={What you've built}]
One more word for this chapter.
\end{grammarlens}

\subsection*{Your turn}
\begin{itemize}
\raggedright
  \item \emph{Say it:} \emph{drākṣi}
  \item \emph{Say it:} \emph{drākṣi}, once more
  \item \emph{Say it:} say \emph{drākṣi}
  \item \emph{Recall:} say the Kannada for finger millet, then the Kannada for a lemon, then say \emph{drākṣi} again
\end{itemize}

\begin{tcolorbox}[breakable,colback=teal!4,colframe=teal!35!black,title={Before you move on}]
What does \kn{ದ್ರಾಕ್ಷಿ} mean? (\textquotedblleft{}Grapes\textquotedblright{}.) Say it once more.
\end{tcolorbox}
\section[\texorpdfstring{\kn{ಹಲಸು} (halasu)}{ಹಲಸು (halasu)}]{\kn{ಹಲಸು} (halasu) — a jackfruit}

\label{lesson:KA-C126-halasu}



\begin{quote}
Before the new one: say the Kannada for a lemon, then the Kannada for grapes.
\end{quote}

\subsection*{You'll want to know: \kn{ಹಲಸು}}
\textbf{\kn{ಹಲಸು}} — \emph{halasu} — \textquotedblleft{}a jackfruit\textquotedblright{}.

\begin{grammarlens}[title={What you've built}]
One more word for this chapter.
\end{grammarlens}

\subsection*{Your turn}
\begin{itemize}
\raggedright
  \item \emph{Say it:} \emph{halasu}
  \item \emph{Say it:} \emph{halasu}, once more
  \item \emph{Say it:} say \emph{halasu}
  \item \emph{Recall:} say the Kannada for a lemon, then the Kannada for grapes, then say \emph{halasu} again
\end{itemize}

\begin{tcolorbox}[breakable,colback=teal!4,colframe=teal!35!black,title={Before you move on}]
What does \kn{ಹಲಸು} mean? (\textquotedblleft{}A jackfruit\textquotedblright{}.) Say it once more.
\end{tcolorbox}
\section[\texorpdfstring{\kn{ಮೆಣಸು} (menasu)}{ಮೆಣಸು (menasu)}]{\kn{ಮೆಣಸು} (menasu) — pepper, chilli}

\label{lesson:KA-C126-menasu}



\begin{quote}
Before the new one: say the Kannada for grapes, then the Kannada for a jackfruit.
\end{quote}

\subsection*{You'll want to know: \kn{ಮೆಣಸು}}
\textbf{\kn{ಮೆಣಸು}} — \emph{menasu} — \textquotedblleft{}pepper, chilli\textquotedblright{}.

\begin{grammarlens}[title={What you've built}]
One more word for this chapter.
\end{grammarlens}

\subsection*{Your turn}
\begin{itemize}
\raggedright
  \item \emph{Say it:} \emph{menasu}
  \item \emph{Say it:} \emph{menasu}, once more
  \item \emph{Say it:} say \emph{menasu}
  \item \emph{Recall:} say the Kannada for grapes, then the Kannada for a jackfruit, then say \emph{menasu} again
\end{itemize}

\begin{tcolorbox}[breakable,colback=teal!4,colframe=teal!35!black,title={Before you move on}]
What does \kn{ಮೆಣಸು} mean? (\textquotedblleft{}Pepper, chilli\textquotedblright{}.) Say it once more.
\end{tcolorbox}
\section[\texorpdfstring{\kn{ಜೀರಿಗೆ} (jīrige)}{ಜೀರಿಗೆ (jīrige)}]{\kn{ಜೀರಿಗೆ} (jīrige) — cumin}

\label{lesson:KA-C126-jirige}



\begin{quote}
Before the new one: say the Kannada for a jackfruit, then the Kannada for pepper.
\end{quote}

\subsection*{You'll want to know: \kn{ಜೀರಿಗೆ}}
\textbf{\kn{ಜೀರಿಗೆ}} — \emph{jīrige} — \textquotedblleft{}cumin\textquotedblright{}.

\begin{grammarlens}[title={What you've built}]
One more word for this chapter.
\end{grammarlens}

\subsection*{Your turn}
\begin{itemize}
\raggedright
  \item \emph{Say it:} \emph{jīrige}
  \item \emph{Say it:} \emph{jīrige}, once more
  \item \emph{Say it:} say \emph{jīrige}
  \item \emph{Recall:} say the Kannada for a jackfruit, then the Kannada for pepper, then say \emph{jīrige} again
\end{itemize}

\begin{tcolorbox}[breakable,colback=teal!4,colframe=teal!35!black,title={Before you move on}]
What does \kn{ಜೀರಿಗೆ} mean? (\textquotedblleft{}Cumin\textquotedblright{}.) Say it once more.
\end{tcolorbox}
\section[(dialogue)]{The house and things you carry — a first pass}

\label{lesson:KA-R126-first-pass-the-house-and-things-you-carry}



\begin{quote}
Say the words for pepper and cumin.
\end{quote}

\subsection*{Your turn}
Say these aloud:

\begin{itemize}
\raggedright
  \item a hammer, scissors, a comb, soap, a towel
  \item glasses, a clock, a bag, a pen, paper; a letter
  \item a stove, a candle, a matchbox, a doll, a ball
  \item a kite, a shirt, a sari, a dhoti, sandals
  \item a cap, a ring, a bangle, rasam, sambar
\end{itemize}

\begin{tcolorbox}[breakable,colback=teal!4,colframe=teal!35!black,title={Before you move on}]
A hammer? (\textbf{\kn{ಸುತ್ತಿಗೆ}}, \emph{suttige}.) Scissors? (\textbf{\kn{ಕತ್ತರಿ}}, \emph{kattari}.) A comb? (\textbf{\kn{ಬಾಚಣಿಗೆ}}, \emph{bācaṇige}.) Soap? (\textbf{\kn{ಸಾಬೂನು}}, \emph{sābūnu}.) A towel? (\textbf{\kn{ಟವೆಲ್}}, \emph{ṭavel}.) Glasses? (\textbf{\kn{ಕನ್ನಡಕ}}, \emph{kannaḍaka}.) A clock? (\textbf{\kn{ಗಡಿಯಾರ}}, \emph{gaḍiyāra}.) A bag? (\textbf{\kn{ಚೀಲ}}, \emph{cīla}.) A pen? (\textbf{\kn{ಲೇಖನಿ}}, \emph{lēkhani}.) Paper; a letter? (\textbf{\kn{ಕಾಗದ}}, \emph{kāgada}.) A stove? (\textbf{\kn{ಒಲೆ}}, \emph{ole}.) A candle? (\textbf{\kn{ಮೇಣಬತ್ತಿ}}, \emph{mēṇabatti}.) A matchbox? (\textbf{\kn{ಬೆಂಕಿಪೆಟ್ಟಿಗೆ}}, \emph{beṅkipeṭṭige}.) A doll? (\textbf{\kn{ಬೊಂಬೆ}}, \emph{bombe}.) A ball? (\textbf{\kn{ಚೆಂಡು}}, \emph{ceṇḍu}.) A kite? (\textbf{\kn{ಗಾಳಿಪಟ}}, \emph{gāḷipaṭa}.) A shirt? (\textbf{\kn{ಅಂಗಿ}}, \emph{aṅgi}.) A sari? (\textbf{\kn{ಸೀರೆ}}, \emph{sīre}.) A dhoti? (\textbf{\kn{ಪಂಚೆ}}, \emph{paṁce}.) Sandals? (\textbf{\kn{ಚಪ್ಪಲಿ}}, \emph{cappali}.) A cap? (\textbf{\kn{ಟೋಪಿ}}, \emph{ṭōpi}.) A ring? (\textbf{\kn{ಉಂಗುರ}}, \emph{uṅgura}.) A bangle? (\textbf{\kn{ಬಳೆ}}, \emph{baḷe}.) Rasam? (\textbf{\kn{ಸಾರು}}, \emph{sāru}.) Sambar? (\textbf{\kn{ಸಾಂಬಾರು}}, \emph{sāmbāru}.)
\end{tcolorbox}
\section[(dialogue)]{Clothes and food — a second pass}

\label{lesson:KA-R126-second-pass-clothes-and-food}



\begin{quote}
Say the words for a dry vegetable dish and a pickle.
\end{quote}

\subsection*{Your turn}
Say these aloud:

\begin{itemize}
\raggedright
  \item a dry vegetable dish, a pickle, an egg, a fish, meat
  \item vegetables, a banana, a mango, a coconut, an onion
  \item a potato, an aubergine, ginger, garlic, honey
  \item an idli, a dosa, a vada, upma, finger millet
  \item a lemon, grapes, a jackfruit, pepper, cumin
\end{itemize}

\begin{tcolorbox}[breakable,colback=teal!4,colframe=teal!35!black,title={Before you move on}]
A dry vegetable dish? (\textbf{\kn{ಪಲ್ಯ}}, \emph{palya}.) A pickle? (\textbf{\kn{ಉಪ್ಪಿನಕಾಯಿ}}, \emph{uppinakāyi}.) An egg? (\textbf{\kn{ಮೊಟ್ಟೆ}}, \emph{moṭṭe}.) A fish? (\textbf{\kn{ಮೀನು}}, \emph{mīnu}.) Meat? (\textbf{\kn{ಮಾಂಸ}}, \emph{māṁsa}.) Vegetables? (\textbf{\kn{ತರಕಾರಿ}}, \emph{tarakāri}.) A banana? (\textbf{\kn{ಬಾಳೆಹಣ್ಣು}}, \emph{bāḷehaṇṇu}.) A mango? (\textbf{\kn{ಮಾವಿನಹಣ್ಣು}}, \emph{māvinahaṇṇu}.) A coconut? (\textbf{\kn{ತೆಂಗು}}, \emph{teṅgu}.) An onion? (\textbf{\kn{ಈರುಳ್ಳಿ}}, \emph{īruḷḷi}.) A potato? (\textbf{\kn{ಆಲೂಗಡ್ಡೆ}}, \emph{ālūgaḍḍe}.) An aubergine? (\textbf{\kn{ಬದನೆಕಾಯಿ}}, \emph{badanekāyi}.) Ginger? (\textbf{\kn{ಶುಂಠಿ}}, \emph{śuṇṭhi}.) Garlic? (\textbf{\kn{ಬೆಳ್ಳುಳ್ಳಿ}}, \emph{beḷḷuḷḷi}.) Honey? (\textbf{\kn{ಜೇನು}}, \emph{jēnu}.) An idli? (\textbf{\kn{ಇಡ್ಲಿ}}, \emph{iḍli}.) A dosa? (\textbf{\kn{ದೋಸೆ}}, \emph{dōse}.) A vada? (\textbf{\kn{ವಡೆ}}, \emph{vaḍe}.) Upma? (\textbf{\kn{ಉಪ್ಪಿಟ್ಟು}}, \emph{uppiṭṭu}.) Finger millet? (\textbf{\kn{ರಾಗಿ}}, \emph{rāgi}.) A lemon? (\textbf{\kn{ನಿಂಬೆ}}, \emph{nimbe}.) Grapes? (\textbf{\kn{ದ್ರಾಕ್ಷಿ}}, \emph{drākṣi}.) A jackfruit? (\textbf{\kn{ಹಲಸು}}, \emph{halasu}.) Pepper? (\textbf{\kn{ಮೆಣಸು}}, \emph{menasu}.) Cumin? (\textbf{\kn{ಜೀರಿಗೆ}}, \emph{jīrige}.)
\end{tcolorbox}
